\section{Day 8: Gifts for the Common Good and the Mission}

\dayaxis{1 Corinthians 12:1--31; Acts 13:1--4; 1 Peter 4:7--11}{Charisms ordered by charity, apostolic faith, discernment, and service}{Offer a concrete skill or hour of work to a parish, neighbor, family, or work of mercy without claiming spiritual status.}

\subsection{Read: 1 Corinthians 12:1--31}

Paul's governing sentence is easily lost amid fascination with gifts: “To each is given the manifestation of the Spirit for the common good” (1 Cor 12:7). The one Spirit distributes; the recipients do not own him. The one Body has many members; an eye cannot dismiss a hand, and the apparently weaker members receive special honor. Charisms disclose dependence.

The list is intentionally varied: wisdom, knowledge, faith in a charismatic mode, healings, mighty works, prophecy, discernment, tongues, interpretation, assistance, administration, and more. Other New Testament lists differ (Rom 12:3--8; Eph 4:7--16; 1 Pet 4:10--11). No list is a personality inventory or exhaustive rank. Ordinary capacities can be taken up by grace for ecclesial service; extraordinary phenomena require sober discernment.

Paul immediately places the “more excellent way” of charity between chapters 12 and 14. A charism is not proof of sanctity, doctrinal impeccability, office, or superiority. Prophecy is tested; tongues are ordered; intelligibility and peace serve the assembly; apostolic teaching governs. The Catechism likewise says charisms are to be received with gratitude and discerned by the Church's pastors (CCC 799--801).

\subsection{Meditation: sent, not self-appointed}

In Acts 13 the Church at Antioch worships and fasts; the Spirit sets apart Barnabas and Saul; the community fasts, prays, lays hands, and sends them. Divine initiative, ecclesial recognition, ascetic readiness, and concrete mission remain together. A private certainty does not erase formation, competence, safeguarding, or lawful commissioning.

Some readers undervalue their gifts because they are quiet: administration, maintenance, hospitality, translation, accounting, caregiving, teaching one child, patient intercession, or the fidelity of chronic illness offered in union with Christ. Others inflate a gift into an identity that cannot accept correction. Both errors make the gift about the possessor.

Ask three questions of a proposed service: Is it true and morally good? Does it meet a real need within my state and competence? Can it be offered under appropriate accountability? If yes, begin modestly. Mission usually grows through fidelity before scale.

Pray for vocations to priesthood, diaconate, consecrated life, marriage, dedicated single life, catechesis, scholarship, works of mercy, and public service. These are not interchangeable, yet the same Spirit orders them toward Christ's Body and the world's salvation.

\novenaexamination{Do I envy a visible gift or despise an ordinary one? Have I used “calling” to evade formation, correction, or authority? Who actually benefits from the way I exercise my ability?}

\bilingualprayer{Proper collect for Day 8}{Project-composed Latin prayer and project English translation; not an approved liturgical collect.}{%
Deus, qui unicúique manifestatiónem Spíritus ad commúnem utilitátem tríbuis: revéla nobis quómodo serviámus, ab ináni glória nos custódi, pastóribus tuis donum discretiónis concéde, et omnia charísmata ita órdina, ut corpus Christi ædificétur et Evangélium usque ad fines terræ annuntiétur. Per Christum Dóminum nostrum. Amen.}{%
O God, who givest to each the manifestation of the Spirit for the common good: reveal how we are to serve; keep us from empty glory; grant thy pastors the gift of discernment; and so order every charism that the Body of Christ may be built and the Gospel proclaimed to the ends of the earth. Through Christ our Lord. Amen.}

\prayerinstruction{Continue with \textit{Veni Creator Spiritus}, the Lord's Prayer, and the Lesser Doxology.}
